\addvspace {10\p@ }
\contentsline {figure}{\numberline {1.1}{\ignorespaces Schematic bubble notation for nested Poisson brackets. Popping a bubble produces oriented causal cuts according to the Leibniz rule.}}{9}{figure.caption.47}%
\contentsline {figure}{\numberline {1.2}{\ignorespaces Schematic graphical content of Murua's recursion. The displayed rooted tree has two allowed partitions: \(p_1\), containing only the root-adjacent edges, and \(p_2\), containing all three edges. In general, rooted partitions \(p\in {\rm Pr}(\tau )\) must contain all root-adjacent edges; for non-rooted trees, \(p\in {\rm P}_{s}(\tau )\) must contain all edges adjacent to the chosen semi-root \(s\).}}{11}{figure.caption.57}%
\contentsline {figure}{\numberline {1.3}{\ignorespaces Causality-prescription diagrams for the IR-sensitive master integral. The Magnus/Murua weights assigned to the two oriented orderings cancel the leading \(1/\epsilon ^2\) pole.}}{11}{figure.caption.61}%
\contentsline {figure}{\numberline {1.4}{\ignorespaces Celestial description of a four-dimensional scattering process. A momentum-space \(n\to m\) scattering amplitude in asymptotically flat spacetime is mapped, through a Mellin transform over the external energies, to a conformal correlator on the celestial sphere.}}{14}{figure.caption.70}%
\contentsline {figure}{\numberline {1.5}{\ignorespaces Geometry and contour of nested integrals}}{18}{figure.caption.106}%
\addvspace {10\p@ }
\contentsline {figure}{\numberline {2.1}{\ignorespaces Tree level diagram (left) vs. loop level (one loop) (right) diagram in WEFT}}{29}{figure.caption.129}%
\contentsline {figure}{\numberline {2.2}{\ignorespaces Different scales in the binary inspiral problem. We reproduce the figure from \cite {Porto:2016pyg}.}}{30}{figure.caption.133}%
\contentsline {figure}{\numberline {2.3}{\ignorespaces Diagrams contributing to the gravitational bound sector at 1PN. }}{38}{figure.caption.164}%
\contentsline {figure}{\numberline {2.4}{\ignorespaces Diagrams contributing to the electromagnetic bound sector up to 1PN.}}{39}{figure.caption.171}%
\contentsline {figure}{\numberline {2.5}{\ignorespaces Diagrams contributing to the Proca bound sector up to 1PN.}}{40}{figure.caption.178}%
\contentsline {figure}{\numberline {2.6}{\ignorespaces Scalar-Electromagnetic interaction diagrams which contribute to the bound sector.}}{40}{figure.caption.183}%
\contentsline {figure}{\numberline {2.7}{\ignorespaces 1PN diagrams for the scalar sector with 3 point vertex coming from scalar-graviton interaction.}}{43}{figure.caption.196}%
\contentsline {figure}{\numberline {2.8}{\ignorespaces 1PN scalar diagrams that contribute to the bound sector.}}{44}{figure.caption.199}%
\contentsline {figure}{\numberline {2.9}{\ignorespaces Scalar diagrams coming from the 3-point vertices that contribute at 2 PN order in scalar effective action.}}{45}{figure.caption.203}%
\contentsline {figure}{\numberline {2.10}{\ignorespaces Diagrams contributing to scalar bound sector at 2 PN order coming from worldline vertices. }}{48}{figure.caption.208}%
\contentsline {figure}{\numberline {2.11}{\ignorespaces Diagrams contributing to the bound sector at 1PN order due to the Proca-electromagnetic interaction term.}}{49}{figure.caption.213}%
\contentsline {figure}{\numberline {2.12}{\ignorespaces Self-energy diagram contributing to the effective action for the source term $J$.}}{54}{figure.caption.238}%
\contentsline {figure}{\numberline {2.13}{\ignorespaces Radiative diagrams for the scalar field that contribute up to $N^{(4)}LO$.}}{56}{figure.caption.249}%
\contentsline {figure}{\numberline {2.14}{\ignorespaces Radiative diagrams for electromagnetic field up to $N^{(3)}LO$. }}{61}{figure.caption.277}%
\contentsline {figure}{\numberline {2.15}{\ignorespaces LO radiation coming from the pure gravitational sector.}}{66}{figure.caption.302}%
\contentsline {figure}{\numberline {2.16}{\ignorespaces Radiative diagrams for gravitational fields.}}{68}{figure.caption.310}%
\contentsline {figure}{\numberline {2.17}{\ignorespaces $\alpha (m;r)$ vs $mr$ plot}}{83}{figure.caption.362}%
\contentsline {figure}{\numberline {2.18}{\ignorespaces $\beta (m;r)$ vs $mr$ plot}}{83}{figure.caption.364}%
\contentsline {figure}{\numberline {2.19}{\ignorespaces Correspondence between PN diagrams (left) and one-loop scalar Feynman diagram (right).}}{84}{figure.caption.365}%
\addvspace {10\p@ }
\contentsline {figure}{\numberline {3.1}{\ignorespaces Sketch of the celestial sphere under consideration.}}{160}{figure.caption.803}%
\contentsline {figure}{\numberline {3.2}{\ignorespaces Figure showing dependency of the integral $I(b)$ as a function of $|b|$ for $m=1\,.$}}{165}{figure.caption.842}%
\addvspace {10\p@ }
\addvspace {10\p@ }
\contentsline {figure}{\numberline {5.1}{\ignorespaces Diagrammatic representation of Lippmann-Schwinger equation}}{250}{figure.caption.1181}%
\addvspace {10\p@ }
\contentsline {figure}{\numberline {6.1}{\ignorespaces Figure describing the $s,t,u$-channel diagram respectively, relevant for our computation }}{263}{figure.caption.1243}%
\contentsline {figure}{\numberline {6.2}{\ignorespaces Figure depicting the chosen contour for dispersion relation in the complex-$\omega $ plane. The cross signs depict the location of the poles.}}{275}{figure.caption.1297}%
\contentsline {figure}{\numberline {6.3}{\ignorespaces Plot depicting the matching for extraction of OPE coefficient using BC expansion (\textbf {red}) and OPE inversion for spin-0 (\textbf {blue}). We have set the conformal dimension for the external primaries to be: $\Delta _{\mathcal {O}}=1+i$.}}{287}{figure.caption.1358}%
\addvspace {10\p@ }
